\documentclass{article}
\usepackage{amsmath,amssymb}
\usepackage{xurl}
\usepackage[hidelinks]{hyperref}
\setlength{\emergencystretch}{2em}
% Shared commands for notes in docs/paper-gaps/.

\DeclareUnicodeCharacter{2080}{\ifmmode{}_{0}\else\textsubscript{0}\fi}
\DeclareUnicodeCharacter{2081}{\ifmmode{}_{1}\else\textsubscript{1}\fi}
\DeclareUnicodeCharacter{2082}{\ifmmode{}_{2}\else\textsubscript{2}\fi}
\DeclareUnicodeCharacter{2099}{\ifmmode{}_{n}\else\textsubscript{n}\fi}

\newcommand{\Lean}{\textsc{Lean}}
\ExplSyntaxOn
\NewDocumentCommand{\leanid}{m}{
  \ifmmode
    \textsf{\detokenize{#1}}
  \else
    \group_begin:
    \urlstyle{sf}
    \tl_set:Nn \l_tmpa_tl {#1}
    \tl_replace_all:Nnn \l_tmpa_tl {\_} {_}
    \exp_args:NV \path \l_tmpa_tl
    \group_end:
  \fi
}
\ExplSyntaxOff
\providecommand{\tr}{\operatorname{tr}}
\newcommand{\tnleanIssueBase}{https://github.com/LionSR/TNLean/issues/}
\newcommand{\tnleanPrBase}{https://github.com/LionSR/TNLean/pull/}
\newcommand{\tnleanDocBase}{https://sirui-lu.com/TNLean/docs/}
\newcommand{\ghissue}[1]{\href{\tnleanIssueBase#1}{GitHub issue \##1}}
\newcommand{\ghpr}[1]{\href{\tnleanPrBase#1}{PR \##1}}
\newcommand{\doclink}[1]{\href{\tnleanDocBase#1.html}{\path{#1}}}

% Dirac notation and the identity operator, as in blueprint/src/macros/common.tex.
% \providecommand keeps note-local definitions authoritative.
\usepackage{dsfont}
\providecommand{\ket}[1]{|#1\rangle}
\providecommand{\bra}[1]{\langle #1|}
\providecommand{\braket}[2]{\langle #1 | #2 \rangle}
\providecommand{\ketbra}[2]{|#1\rangle\!\langle #2|}
\providecommand{\Id}{\mathds{1}}
\providecommand{\one}{\mathds{1}}
\providecommand{\C}{\mathbb{C}}
\providecommand{\MN}[1]{M_{#1}(\C)}
\providecommand{\MD}{M_D(\C)}

% External referencing for paper sources.  The build helper writes sanitized
% auxiliary files under build/paper-aux/ and excludes duplicated source labels.
\usepackage{xr}

% Import a generated paper auxiliary file with a caller-supplied label prefix.
% The candidate paths support builds from docs/paper-gaps/, the repository root,
% and build/paper-aux/ itself.
\newcommand{\importpaperaux}[2]{%
  \IfFileExists{../../build/paper-aux/#2.aux}{%
    \externaldocument[#1]{../../build/paper-aux/#2}%
  }{%
    \IfFileExists{build/paper-aux/#2.aux}{%
      \externaldocument[#1]{build/paper-aux/#2}%
    }{%
      \IfFileExists{#2.aux}{%
        \externaldocument[#1]{#2}%
      }{}%
    }%
  }%
}

% General external reference macros
\newcommand{\paperref}[2]{\ref{#1-#2}}
\newcommand{\papereqref}[2]{(\ref{#1-#2})}

% CPSV16 source (1606.00608 / MPDO-22-12-17-2.tex) external document and shortcuts
\importpaperaux{cpsv16-}{1606.00608/MPDO-22-12-17-2-tnlean}
\newcommand{\cpsvref}[1]{\paperref{cpsv16}{#1}}
\newcommand{\cpsveqref}[1]{\papereqref{cpsv16}{#1}}

% Verdict marker: \gapnote{<kind>}{<status>}. Kind is the policy
% classification (clarification, local-correction, scope-restriction,
% unfaithful, false-source, open-gap); status is open, wip (elimination
% underway), resolved, or historical. Machine-read by the site index;
% prints a verdict line.
\newcommand{\gapnote}[2]{\par\noindent\textbf{Verdict:} #1 (#2).\par}

\providecommand{\gapnote}[2]{%
    \par\noindent\textbf{Verdict:} #1 (#2).\par}

\title{The Interaction Range in the Block-Injective Parent-Hamiltonian Gap}
\author{}
\date{2026-09-26}

\begin{document}
\maketitle
\gapnote{scope-restriction}{resolved}

\section{The range condition}

The gap theorem in Ref.~\cite[Section~IV.C, Gaps]{Cirac2021Matrix},
local source lines~2183--2187, asserts a uniform positive gap for MPS
parent Hamiltonians, including those of block-injective tensors.
For at least two canonical blocks, Ref.~\cite[Theorem~12]{PerezGarcia2007MPSReps}
gives the explicit interaction condition
\begin{align}
    R\geq 3(r-1)(L_0+1)+1,
    \label{eq:block_range_source_bound}
\end{align}
where the blocks satisfy condition C1 at the common positive length $L_0$.
The local source is \texttt{MPSarchive.tex}, lines~1424--1458.

Let $A^0,\ldots,A^{r-1}$ have positive bond dimensions and be pairwise
inequivalent under virtual conjugation and nonzero scalar multiplication.
For nonzero coefficients $\mu_j$, write $B=\bigoplus_j\mu_j A^j$ and let
$H_{R,N}$ be its canonical periodic parent Hamiltonian.
The overlap argument initially constructs a gapped range $R=2p$ after
choosing $p$ sufficiently large. This alone would not prove a gap at every
prescribed range satisfying (\ref{eq:block_range_source_bound}).
The comparison below removes this restriction, as developed under
\href{https://github.com/LionSR/TNLean/issues/7676}{issue \#7676}.

\section{Comparison of two interaction ranges}

Fix $R>0$, and choose $W\geq 2R$ such that the open range-$R$ Hamiltonian
$K_{R,W}$ has kernel $G_W(B)$ and the periodic range-$W$ Hamiltonians
have a uniform gap $\gamma_W>0$ for all sufficiently large volumes.
Set $h_W=I-P_{G_W(B)}$. Finite dimensionality and positivity give constants
$\kappa,C>0$ such that
\begin{align}
    \kappa h_W\leq K_{R,W}\leq C h_W.
    \label{eq:block_range_local_comparison}
\end{align}
The lower bound follows from injectivity of $K_{R,W}$ on its kernel
complement. For the upper bound one may take
$C=\lVert K_{R,W}\rVert+1$.

Embed (\ref{eq:block_range_local_comparison}) at every cyclic initial site
of a sufficiently long periodic chain. Each range-$R$ term appears
exactly $m=W-R+1$ times. Thus
\begin{align}
    \kappa H_{W,N}\leq mH_{R,N}\leq C H_{W,N}.
    \label{eq:block_range_periodic_comparison}
\end{align}
Both periodic operators have the same kernel: the quadratic form of a
positive operator vanishes precisely on its kernel, and the two inequalities
give the two inclusions. The range-$W$ gap therefore implies
\begin{align}
    \frac{\kappa\gamma_W}{m}\lVert v\rVert
    \leq\lVert H_{R,N}v\rVert
    \qquad(v\perp\ker H_{R,N}).
    \label{eq:block_range_transferred_gap}
\end{align}
All constants are independent of $N$. To include the finitely many excluded
lengths, restrict each corresponding Hamiltonian to its kernel complement.
This restriction is injective and finite-dimensional, so it has a positive
norm lower bound. Taking the minimum of these finitely many bounds and the
uniform eventual bound gives one positive constant for every $N$. The
comparison uses every cyclic
initial site and imposes no divisibility restriction on the chain length.
It never identifies the open-chain kernel with the periodic kernel.

\section{The simultaneous-injectivity range}

Suppose instead that a positive integer $S$ is supplied for which the length-$S$
word tuples of the blocks span the full product matrix algebra. This is the
simultaneous block-injectivity condition in Ref.~\cite[Section~IV.C]{Cirac2021Matrix},
local source lines~2114--2129. Then the uniform gap holds at every range
\begin{align}
    R\geq S+1.
    \label{eq:block_range_simultaneous_injectivity}
\end{align}
The number of blocks need not be at least two, and no normalization or separate
inequivalence assumption is imposed on the original blocks.

Indeed, projecting the simultaneous span onto any component proves injectivity
of that block at length $S$, so each periodic vector $V^{(S)}(A^i)$ is
nonzero. The simultaneous span also makes the local ground spaces $G_S(A^i)$
independent, and this separates distinct blocks. Suppose that $i\neq j$ and
$A^i_a=\zeta XA^j_aX^{-1}$ for every physical index $a$, with $\zeta\neq0$ and
$X$ invertible. Then
\begin{align}
    V^{(S)}(A^i)=\zeta^SV^{(S)}(A^j)\in G_S(A^i)\cap G_S(A^j)=\{0\},
    \label{eq:block_range_separation}
\end{align}
contradicting $V^{(S)}(A^i)\neq0$. Independent
primitive normalization preserves the supplied length $S$, simultaneous
spanning, and all local ground spaces. In the normalized family, simultaneous
spanning propagates to every larger length. The open-chain intersection theorem
therefore gives $\ker K_{R,W}=G_W(B)$ for every $W\geq R$ as soon as (\ref{eq:block_range_simultaneous_injectivity}) holds.
Choose a gapped range $W\geq2R$ and apply the two-sided comparison above.
The finite minimum over the remaining chain lengths again gives one positive
gap for every $N$.

For an original canonical tensor with repeated multiplicity copies, the local
ground spaces equal those of its chosen normal representatives at every
positive length. Hence the canonical parent interactions, and then every
periodic parent Hamiltonian, agree exactly. The same uniform gap transfers to
the original tensor. The representatives may be chosen once, before specifying
$S$ or $R$.

\section{General positive parent interactions}

The definition in Ref.~\cite[Section~IV.C]{Cirac2021Matrix}, local source
lines~1996--1999, permits any fixed positive local operator $h$ with
$\ker h=G_R(A)$. Let $q_R$ be the orthogonal projection onto $G_R(A)^\perp$.
Finite dimensionality gives constants $\kappa,C>0$ such that
\begin{align}
    \kappa q_R\leq h\leq Cq_R.
\end{align}
If $h=0$, then $q_R=0$ and any positive constants suffice. Otherwise one may
use the least and greatest positive eigenvalues of $h$.

Tensoring with the identity, conjugating by cyclic translations, and summing
over initial sites preserve these inequalities. Thus, with the same constants
for every chain length $N\geq R$,
\begin{align}
    \kappa H_{R,N}(A)\leq H_N(h)\leq C H_{R,N}(A).
\end{align}
The two Hamiltonians have the same kernel. A uniform gap $\gamma$ for the
canonical projection Hamiltonian therefore gives a uniform gap
$\kappa\gamma$ for $H_N(h)$. This is the projection reduction stated in the
review at local source lines~2170--2172. In particular, it preserves the
interaction range $R\geq S+1$ established above.

The interaction $h$ is fixed before varying $N$. Equality of local kernels
alone would not suffice for an arbitrary volume-dependent choice: replacing
$h$ by $h/N$ preserves each kernel but need not preserve a positive uniform gap.

\section{Recovering the printed threshold}

The simultaneous word tuples span the full product matrix algebra at length
$3(r-1)(L_0+1)$. Trace-preserving normalization propagates this span to every
larger length: represent $M_j(A^j_a)^\dagger$ by word tuples, append $A^j_a$,
and sum over $a$, using $\sum_a(A^j_a)^\dagger A^j_a=I$.
The intersection property then identifies the open range-$R$ kernel with
$G_W(B)$ for every $W\geq R$ whenever (\ref{eq:block_range_source_bound}) holds.
The overlap construction supplies a gapped $W$ above any prescribed bound,
so the preceding comparison applies to every such $R$.

Finally, independent primitive normalization of the blocks preserves
injectivity at the specified length $L_0$, pairwise inequivalence, and all
local ground spaces. It therefore preserves every parent Hamiltonian exactly.
No larger C1 length is substituted during this step. The formalized conclusion
is consequently a uniform periodic gap over every chain length at each range in
(\ref{eq:block_range_source_bound}), with no normalization hypothesis on the
original blocks. The earlier chosen-range restriction is resolved.

\bibliographystyle{alpha}
\bibliography{references}
\end{document}
